\chapter*{Conclusion \& Perspectives \markboth{Conclusion}{}}
\addcontentsline{toc}{chapter}{Conclusion}
\setcounter{section}{0}

This thesis work has proposed the development and validation of several numerical tools allowing to model metaporous surfaces and interfaces. These structures feature dissipation mechanisms, strong dispersion due to the periodicity as well as leakage, thus involving challenges in the modeling process. Methods were developed to compute the dispersion relations of leaky guided waves propagating in homogeneous multilayer structures, evaluate the scattering properties of grating stacks and compute dispersion relations of waves in metaporous layers.

Throughout the literature, these metaporous structures have shown efficient absorption properties using various geometry by embedding resonant, tortuous paths, voids or inclusions inside porous and poroelastic layers. These obstacles provide additional dispersion as well as interferences of the wave at discrete frequencies. Combined with the attenuation of waves in porous layer, they provide for efficient acoustic absorption structures. Even though, these effects have been studied in the literature, they have mostly been considered as infinite and closed systems. In this work, the considerations were extended to the study of leaky guided waves radiating to surrounding media. The combination of these loss mechanisms thus have showned a complex behaviour featured by the structure and that the classification of the waves was of interest and that the physical mechanisms at stake were not so simple. 

The foundation for the work was laid out in \textsc{chapter} 1, introducing the governing motion equations for the assumptions of homogeneous porous and poroelastic media and proposed a state of art for the exisitng numerical method allowing for the computation of the scattering properties and the dispersion of waves in homogeneous layers and periodic systems. 

Next, \textsc{chapter 2} deals with dissipative homogeneous per layer multilayer structures accounting for leakage. These homogeneous layers can be described by prescribing a harmonic motion in the infinite direction while using the SCM and a 1D discretization to describe the propagation along the transverse direction. With this general formalism, any stack of layer, dissipative or not, can be computed with the presented method. Results regarding the energy flux for specific modes were presented as well. This method was implemented and validated as a fast and accurate technique for reliably computing the dispersion relations. Its efficiency and stability were demonstrated through comparisons with semi-analytical results. Results include the study of a two-layer structure\comment{finir ça.}

In \textsc{chapter} 3, a robust and stable method for dealing with grating stacks constituted of poroelastic and porous layers. finite element discretization remained stable for all tested cases contrary to the TMM and the coupling is thus quite robust and efficient, since no fem discretization of the homogeneous layers is required. 

In the final \textsc{chapter 4} we adapted the previous method, in the absence of a source, solving its dispersion relation. The method relies on the same discretization as previously, therefore the root-finding solving of the method is not very efficient. However, reliable results have still been achieved. We were able to highlight the trapped mode that occurs in a perfect absorption case as well as the impact of the resonance due to an hybrization of the resonator and the cavity for split-ring resonators. 


\subsection*{Perspectives and future work}

\textbf{1.} The finite elements approach developed and presented in Chapter 3 has been tested and validated for general layers. However, in Chapter 4, dispersion relations are only studied for the case of equivalent fluid; future work thus include an extension of the study of dispersion relation of waves in metaporous structures, by assuming poroelastic layers and solid inclusions, leading to a much larger of modes to analyse due to the presence of the modes linked to the elastic phase of the materials modeled with the full Biot theory. As observed for homogeneous layers, this leads to a much larger number of modes to be solved. While the numerical implementation does not consistute an issue as the basis of it was validated in Chapter 3, the analysis of the results is expected to be quite complex.
\comment{investigate more deeply the physics, energy conversion mechanisms and such}

\textbf{2.} An important perspective for the metaporous structure dispersion relation computation presented in Chapter 4 would involve the development of a more robust method not relying on a root-finding procedure. Some time has been spent during the development of the method to try to write the equations in the form of an eigenvalue problem with $k_{1q}$ as the eigenvalue , although this was not achieve due to the challenges of obtaining such a formulation with the way the coupling between the surrounding medium and the structure is done. The finite elements matrices obtained in Eq.~\comment{ref the correct eq} depend on both on a term in $ e^{ik_{1q}d} $ and an integral on $ x_1 $  with a $ e^{ik_{1q}}x_1 $ term. An appropriate change of variable or a different formulation for the coupling might provide a way to achieve the writing of an eigenvalue problem. Note that this has been explained briefly in Section.~\comment{ref annexe PRB} and shows that the eigenvalue problem is easily written when not accounting for the surrounding medium. 

Methods like contour integral evaluation by using the Beyn method or the FEAST eigenvalue solver. \cite{guttel2017, liu2025} Although these iterative procedures have the advantage of not requiring initial guesses as initialisation, the contour integrals are still computed around a given value, thus still requiring some sort of initialization. Additionally, the quality of the matrix conditioning should be investigated as well due to the amount of sensitive operations required by these methods (e.g. the successive SVD should be computed with sufficient precision). 

\textbf{3.} A study the excitability of the modes could be conducted by using a numerical model done with COMSOL \& a line source excitation properly tuned according to the dispersion relation modes to confirm the form of the solutions of the modes in finite structures and a more \emph{realistic} configuration and further link the dispersion relation results to an actual observable behaviour of the structures by gaining insights on the excitability of the modes. This is an important aspect as there are many dissipation mechanisms involved, leading to very attenuated modes: seeing if they can be excited or if their existence actually affects the wave propagation in any way. 

\comment{near field nature of leaky modes, wannier functions, extension to 3d w/ out of plane modes} Simulation results in these situations could allow for further experimental validation for the configurations of interest. One could imagine a porous sample manufactured with rigid cylindrical inclusions. There would be challenges to this, including the precision required to engineer precisely the radius of the cylinders. Measurements need to take into account the nature of the source needed to excite the variety of these modes. A post-process of the frequency responses using SLaTCoW or some other method allowing for the retrieval of complex wavenumbers. This approach is similar to that applied in Section.\ref{sec:JAP_slatcow}.